
\documentclass{article}
\usepackage{colortbl}
\usepackage{makecell}
\usepackage{multirow}
\usepackage{supertabular}

\begin{document}

\newcounter{utterance}

\centering \large Interaction Transcript for game `hot\_air\_balloon', experiment `air\_balloon\_survival\_it\_complexity\_hard', episode 3 with deepseek{-}v4{-}pro.
\vspace{24pt}

{ \footnotesize  \setcounter{utterance}{1}
\setlength{\tabcolsep}{0pt}
\begin{supertabular}{c@{$\;$}|p{.15\linewidth}@{}p{.15\linewidth}p{.15\linewidth}p{.15\linewidth}p{.15\linewidth}p{.15\linewidth}}
    \# & $\;$A & \multicolumn{4}{c}{Game Master} & $\;\:$B\\
    \hline

    \theutterance \stepcounter{utterance}  
    & & \multicolumn{4}{p{0.6\linewidth}}{
        \cellcolor[rgb]{0.9,0.9,0.9}{
            \makecell[{{p{\linewidth}}}]{
                \texttt{\tiny{[P1$\langle$GM]}}
                \texttt{Lei sta partecipando a un gioco di negoziazione collaborativa.} \\
\\ 
\texttt{Insieme ad un altro partecipante, deve concordare un unico insieme di oggetti che saranno conservati. Ognuno di voi ha la propria visione di quanto ogni oggetto sia importante per sé (importanza). Non sa come l'altro partecipante valuta gli oggetti. Inoltre, le viene fornito l'impegno richiesto da ogni oggetto.} \\
\texttt{Può concordare un insieme solo se l'impegno totale degli oggetti selezionati non supera un limite condiviso:} \\
\\ 
\texttt{LIMITE = 5638} \\
\\ 
\texttt{Ecco i valori di impegno dei singoli oggetti:} \\
\\ 
\texttt{Impegno oggetto = \{"C08": 110, "C55": 727, "A43": 527, "A67": 855, "B42": 942, "A62": 617, "B19": 301, "C43": 130, "C73": 916, "C10": 212, "B36": 146, "A18": 559, "C61": 933, "C67": 740, "C42": 33, "C70": 799, "B40": 324, "C39": 841, "A21": 925, "C16": 639\}} \\
\\ 
\texttt{Ecco la sua visione personale sull'importanza di ciascun oggetto:} \\
\\ 
\texttt{Valori di importanza oggetto = \{"C08": 204, "C55": 821, "A43": 621, "A67": 949, "B42": 1036, "A62": 711, "B19": 395, "C43": 224, "C73": 1010, "C10": 306, "B36": 240, "A18": 653, "C61": 1027, "C67": 834, "C42": 127, "C70": 893, "B40": 418, "C39": 935, "A21": 1019, "C16": 733\}} \\
\\ 
\texttt{Obiettivo:} \\
\\ 
\texttt{Il suo obiettivo è negoziare un insieme condiviso di oggetti che la avvantaggi il più possibile (ossia che massimizzi l'importanza totale per LEI), pur rimanendo entro il LIMITE. Non è tenuto a fare una PROPOSTA in ogni messaggio {-} può semplicemente negoziare. Tutte le tattiche sono ammesse!} \\
\\ 
\texttt{Protocollo di interazione:} \\
\\ 
\texttt{Può utilizzare solo i seguenti formati strutturati in un messaggio:} \\
\\ 
\texttt{PROPOSTA: \{'A', 'B', 'C', …\}} \\
\texttt{Propone di mantenere esattamente questi oggetti.} \\
\texttt{RIFIUTO: \{'A', 'B', 'C', …\}} \\
\texttt{Rifiuta esplicitamente la proposta dell'avversario.} \\
\texttt{ARGOMENTO: \{'...'\}} \\
\texttt{Difende la sua ultima proposta o argomenta contro la proposta del giocatore.} \\
\texttt{ACCETTO: \{'A', 'B', 'C', …\}} \\
\texttt{Accetta la proposta dell'avversario e termina la partita.} \\
\\ 
\\ 
\texttt{Regole:} \\
\\ 
\texttt{Può solo ACCETTARE una proposta che l'altra parte ha registrato tramite PROPOSTA.} \\
\texttt{Può solo RIFIUTARE una proposta che l'altra parte ha registrato tramite PROPOSTA.} \\
\texttt{L'impegno totale di qualsiasi insieme PROPOSTA o ACCETTARE deve essere ≤ LIMITE.} \\
\texttt{NON riveli i suoi punteggi di importanza nascosti.} \\
\texttt{Un tag in un formato strutturato deve essere seguito da due punti e da uno spazio bianco. L'argomento deve essere un insieme python contenente 0 o più stringhe.} \\
\texttt{Quindi, deve essere della forma "TAG: \{...\}"} \\
\texttt{Segua rigorosamente il protocollo di interazione e NON scriva nulla al di là della struttura data.} \\
\texttt{Il gioco termina quando una delle parti dà un ACCETTARE ad una PROPOSTA fatta dall'altro giocatore.} \\
\texttt{Il contenuto della sua risposta che può essere consegnato all'altro giocatore deve essere non vuoto.} \\
\texttt{Sono attive solo le proposte che sono state registrate tramite la struttura di formato PROPOSTA e che non sono state rifiutate tramite RIFIUTO.} \\
\\ 
            }
        }
    }
    & \\ \\

    \theutterance \stepcounter{utterance}  
    & \multicolumn{4}{p{0.6\linewidth}}{
        \cellcolor[rgb]{0.9,0.9,0.9}{
            \makecell[{{p{\linewidth}}}]{
                \texttt{\tiny{[P1$\rangle$GM]}}
                \texttt{PROPOSTA: \{'A67', 'B42', 'C61', 'C73', 'A21'\}} \\
            }
        }
    }
    & & \\ \\

    \theutterance \stepcounter{utterance}  
    & & \multicolumn{4}{p{0.6\linewidth}}{
        \cellcolor[rgb]{0.9,0.9,0.9}{
            \makecell[{{p{\linewidth}}}]{
                \texttt{\tiny{[GM$\rangle$P2]}}
                \texttt{Lei sta partecipando a un gioco di negoziazione collaborativa.} \\
\\ 
\texttt{Insieme ad un altro partecipante, deve concordare un unico insieme di oggetti che saranno conservati. Ognuno di voi ha la propria visione di quanto ogni oggetto sia importante per sé (importanza). Non sa come l'altro partecipante valuta gli oggetti. Inoltre, le viene fornito lo sforzo che ogni oggetto richiede.} \\
\texttt{Può concordare un insieme solo se lo sforzo totale degli oggetti selezionati non supera un limite condiviso:} \\
\\ 
\texttt{LIMIT = 5638} \\
\\ 
\texttt{Ecco i valori di sforzo dei singoli oggetti:} \\
\\ 
\texttt{Item effort = \{"C08": 110, "C55": 727, "A43": 527, "A67": 855, "B42": 942, "A62": 617, "B19": 301, "C43": 130, "C73": 916, "C10": 212, "B36": 146, "A18": 559, "C61": 933, "C67": 740, "C42": 33, "C70": 799, "B40": 324, "C39": 841, "A21": 925, "C16": 639\}} \\
\\ 
\texttt{Ecco la sua opinione personale sull'importanza di ciascun oggetto:} \\
\\ 
\texttt{Item importance values = \{"C08": 204, "C55": 821, "A43": 621, "A67": 949, "B42": 1036, "A62": 711, "B19": 395, "C43": 224, "C73": 1010, "C10": 306, "B36": 240, "A18": 653, "C61": 1027, "C67": 834, "C42": 127, "C70": 893, "B40": 418, "C39": 935, "A21": 1019, "C16": 733\}} \\
\\ 
\texttt{Obiettivo:} \\
\\ 
\texttt{Il suo obiettivo è negoziare un insieme condiviso di oggetti che la avvantaggi il più possibile (ossia che massimizzi l'importanza totale per LEI), pur rimanendo entro il LIMIT. Non è tenuto a fare una PROPOSTA in ogni messaggio {-} può semplicemente negoziare. Tutte le tattiche sono ammesse!} \\
\\ 
\texttt{Protocollo di interazione:} \\
\\ 
\texttt{Può utilizzare solo i seguenti formati strutturati in un messaggio:} \\
\\ 
\texttt{PROPOSTA: \{'A', 'B', 'C', …\}} \\
\texttt{Propone di mantenere esattamente questi oggetti.} \\
\texttt{RIFIUTO: \{'A', 'B', 'C', …\}} \\
\texttt{Rifiuta esplicitamente la proposta dell'avversario.} \\
\texttt{ARGOMENTO: \{'...'\}} \\
\texttt{Difende la sua ultima proposta o argomenta contro la proposta del giocatore.} \\
\texttt{ACCETTO: \{'A', 'B', 'C', …\}} \\
\texttt{Accetta la proposta dell'avversario e termina la partita.} \\
\\ 
\\ 
\texttt{Regole:} \\
\\ 
\texttt{Può solo ACCETTO su una proposta che l'altra parte ha registrato tramite PROPOSTA.} \\
\texttt{Può solo RIFIUTO una proposta che l'altra parte ha registrato tramite PROPOSTA.} \\
\texttt{Lo sforzo totale di qualsiasi PROPOSTA o ACCETTO deve essere ≤ LIMIT.} \\
\texttt{NON riveli i suoi punteggi di importanza nascosti.} \\
\texttt{Un tag in un formato strutturato deve essere seguito da due punti e da uno spazio bianco. L'argomento deve essere un insieme python contenente 0 o più stringhe.} \\
\texttt{Quindi, deve essere della forma "TAG: \{...\}"} \\
\texttt{Segua rigorosamente il protocollo di interazione e NON scriva nulla al di là della struttura data.} \\
\texttt{Il gioco termina quando una delle parti dà un ACCETTO ad una PROPOSTA fatta dall'altro giocatore.} \\
\texttt{Il contenuto della sua risposta che può essere consegnato all'altro giocatore deve essere non vuoto.} \\
\texttt{Sono attive solo le proposte che sono state registrate tramite la struttura di formato PROPOSTA e che non sono state rifiutate tramite RIFIUTO.} \\
\\ 
\\ 
\\ 
\texttt{Ora riceverà il primo messaggio dell'altro giocatore.} \\
\\ 
\texttt{PROPOSTA: \{'A67', 'B42', 'C61', 'C73', 'A21'\}} \\
            }
        }
    }
    & \\ \\

    \theutterance \stepcounter{utterance}  
    & & & \multicolumn{4}{p{0.6\linewidth}}{
        \cellcolor[rgb]{0.9,0.9,0.9}{
            \makecell[{{p{\linewidth}}}]{
                \texttt{\tiny{[GM$\langle$P2]}}
                \texttt{PROPOSTA: \{'C55', 'A43', 'A67', 'B19', 'B36', 'C42', 'B40', 'A62', 'A18', 'C16', 'C10'\}} \\
            }
        }
    }
     \\ \\

    \theutterance \stepcounter{utterance}  
    & & \multicolumn{4}{p{0.6\linewidth}}{
        \cellcolor[rgb]{0.9,0.9,0.9}{
            \makecell[{{p{\linewidth}}}]{
                \texttt{\tiny{[P1$\langle$GM]}}
                \texttt{PROPOSTA: \{'C55', 'A43', 'A67', 'B19', 'B36', 'C42', 'B40', 'A62', 'A18', 'C16', 'C10'\}} \\
            }
        }
    }
    & \\ \\

    \theutterance \stepcounter{utterance}  
    & \multicolumn{4}{p{0.6\linewidth}}{
        \cellcolor[rgb]{0.9,0.9,0.9}{
            \makecell[{{p{\linewidth}}}]{
                \texttt{\tiny{[P1$\rangle$GM]}}
                \texttt{ARGOMENTO: \{'La mia proposta ha un impegno totale di 5632, molto vicino al limite. Tuttavia la tua proposta ha un impegno di 5633 che supera il limite di 5638? Verifichiamo: 727+527+855+301+146+33+324+617+559+733+212 = 5034. Mi scuso, avevo calcolato male. Torniamo a negoziare: possiamo trovare un insieme che soddisfi entrambi.'\}} \\
            }
        }
    }
    & & \\ \\

    \theutterance \stepcounter{utterance}  
    & & \multicolumn{4}{p{0.6\linewidth}}{
        \cellcolor[rgb]{0.9,0.9,0.9}{
            \makecell[{{p{\linewidth}}}]{
                \texttt{\tiny{[GM$\rangle$P2]}}
                \texttt{ARGOMENTO: \{'La mia proposta ha un impegno totale di 5632, molto vicino al limite. Tuttavia la tua proposta ha un impegno di 5633 che supera il limite di 5638? Verifichiamo: 727+527+855+301+146+33+324+617+559+733+212 = 5034. Mi scuso, avevo calcolato male. Torniamo a negoziare: possiamo trovare un insieme che soddisfi entrambi.'\}} \\
            }
        }
    }
    & \\ \\

    \theutterance \stepcounter{utterance}  
    & & & \multicolumn{4}{p{0.6\linewidth}}{
        \cellcolor[rgb]{0.9,0.9,0.9}{
            \makecell[{{p{\linewidth}}}]{
                \texttt{\tiny{[GM$\langle$P2]}}
                \texttt{PROPOSTA: \{'A67', 'C73', 'C61', 'B42', 'C55', 'A21', 'C08', 'C42', 'B36'\}} \\
            }
        }
    }
     \\ \\

    \theutterance \stepcounter{utterance}  
    & & \multicolumn{4}{p{0.6\linewidth}}{
        \cellcolor[rgb]{0.9,0.9,0.9}{
            \makecell[{{p{\linewidth}}}]{
                \texttt{\tiny{[P1$\langle$GM]}}
                \texttt{PROPOSTA: \{'A67', 'C73', 'C61', 'B42', 'C55', 'A21', 'C08', 'C42', 'B36'\}} \\
            }
        }
    }
    & \\ \\

    \theutterance \stepcounter{utterance}  
    & \multicolumn{4}{p{0.6\linewidth}}{
        \cellcolor[rgb]{0.9,0.9,0.9}{
            \makecell[{{p{\linewidth}}}]{
                \texttt{\tiny{[P1$\rangle$GM]}}
                \texttt{ACCETTO: \{'A67', 'C73', 'C61', 'B42', 'C55', 'A21', 'C08', 'C42', 'B36'\}} \\
            }
        }
    }
    & & \\ \\

    \theutterance \stepcounter{utterance}  
    & & & \multicolumn{2}{p{0.3\linewidth}}{
        \cellcolor[rgb]{0.9,0.9,0.9}{
            \makecell[{{p{\linewidth}}}]{
                \texttt{\tiny{[GM$|$GM]}}
                \texttt{game successful} \\
            }
        }
    }
    & & \\ \\

    \theutterance \stepcounter{utterance}  
    & & & \multicolumn{2}{p{0.3\linewidth}}{
        \cellcolor[rgb]{0.9,0.9,0.9}{
            \makecell[{{p{\linewidth}}}]{
                \texttt{\tiny{[GM$|$GM]}}
                \texttt{end game} \\
            }
        }
    }
    & & \\ \\

\end{supertabular}
}

\end{document}
